\documentclass[a4paper]{article}
% generated by Docutils <http://docutils.sourceforge.net/>
\usepackage{cmap} % fix search and cut-and-paste in Acrobat
\usepackage{ifthen}
\usepackage[T1]{fontenc}
\usepackage[utf8]{inputenc}
\setcounter{secnumdepth}{0}

%%% Custom LaTeX preamble
% PDF Standard Fonts
\usepackage{mathptmx} % Times
\usepackage[scaled=.90]{helvet}
\usepackage{courier}

%%% User specified packages and stylesheets

%%% Fallback definitions for Docutils-specific commands

% titlereference role
\providecommand*{\DUroletitlereference}[1]{\textsl{#1}}

% hyperlinks:
\ifthenelse{\isundefined{\hypersetup}}{
  \usepackage[colorlinks=true,linkcolor=blue,urlcolor=blue]{hyperref}
  \usepackage{bookmark}
  \urlstyle{same} % normal text font (alternatives: tt, rm, sf)
}{}
\hypersetup{
  pdftitle={RAT-PAC / ROOT IBD Event Viewer},
}

%%% Body
\begin{document}
\title{RAT-PAC / ROOT IBD Event Viewer%
  \label{rat-pac-root-ibd-event-viewer}}
\author{}
\date{}
\maketitle

Instructions for using Mark Duvall's ROOT-based event viewer for RAT-PAC IBD events.

\textbf{Written for ROOT 5.34; ROOT 6 version currently in development}

\textbf{NOTE: This document presupposes that a working installation of RAT-PAC and its accompanying build of ROOT are active on the machine in use. In other words, the version of ROOT used \_must\_ be one compiled against RAT-PAC, as described in the RAT-PAC documentation {[}here{]}(https://rat.readthedocs.io/en/latest/installation.html).}


\section{Installation%
  \label{installation}%
}

\#\#\# Automatic

For users running a BASH-type shell, there is an installer script on GitHub at \emph{duvall3 -{}-> rat-pac -{}-> branch collab -{}-> user -{}-> shell\_scripts -{}-> {[}RPEV\_Installer.sh{]}(https://github.com/duvall3/rat-pac/blob/collab/user/shell\_scripts/RPEV\_Installer.sh) ({[}file{]}(https://github.com/duvall3/rat-pac/raw/collab/user/shell\_scripts/RPEV\_Installer.sh))}.

\#\#\# Manual

The viewer is a mostly self-contained ROOT macro. It is located on GitHub at \emph{duvall3 -{}-> rat-pac -{}-> branch collab -{}-> user -{}-> root\_macros -{}-> {[}RATPACEventviewer.cxx{]}(https://github.com/duvall3/rat-pac/blob/collab/user/root\_macros/RATPACEventViewer2.cxx) ({[}file{]}(https://github.com/duvall3/rat-pac/raw/collab/user/root\_macros/RATPACEventViewer2.cxx))}.
Remember that the macro (or a link to it) must be located in your \DUroletitlereference{\$ROOTSYS/macros} directory, or you will have to specify the full path to the macro file.
As of the current version, several other files are also rerquired. Their names and locations can be found in the installer script mentioned above.
% These files must also be located (or linked) in \DUroletitlereference{\$ROOTSYS/macros}, as with the main file itself.
% They can be found at the following locations:
% * {[}drawTracks.cxx{]}(\url{https://github.com/duvall3/rat-pac/blob/collab/user/root_macros/drawTracks.cxx}) ({[}file{]}(\url{https://github.com/duvall3/rat-pac/raw/collab/user/root_macros/drawTracks.cxx}))
% * {[}drawNextEvent.cxx{]}(\url{https://github.com/duvall3/rat-pac/blob/collab/user/root_macros/drawNextEvent.cxx}) ({[}file{]}(\url{https://github.com/duvall3/rat-pac/raw/collab/user/root_macros/drawNextEvent.cxx}))
% * {[}drawPrevEvent.cxx{]}(\url{https://github.com/duvall3/rat-pac/blob/collab/user/root_macros/drawPrevEvent.cxx}) ({[}file{]}(\url{https://github.com/duvall3/rat-pac/raw/collab/user/root_macros/drawPrevEvent.cxx}))
% * {[}highlightCell.cxx{]}(\url{https://github.com/duvall3/rat-pac/blob/collab/user/root_macros/highlightCells.cxx}) ({[}file{]}(\url{https://raw.githubusercontent.com/duvall3/rat-pac/collab/user/root_macros/highlightCells.cxx}))


\section{Usage%
  \label{usage}%
}

\#\#\# Basic Invocation

To use the viewer, simply load the macro and run it on the desired ROOT file.  To invoke it in its default configuration, simply type the following at the ROOT prompt:

\texttt{`cpp
.L RATPACEventViewer2.cxx
RATPACEventViewer2(\textquotedbl{}some\_data\_file.root\textquotedbl{});
`}
\#\#\# General Invocation

The viewer decides which detector volumes to draw by looking for a regular expression in the volume names, which by default is \DUroletitlereference{.*target\_cell\_{[}0-9{]}.*}.
This default option matches any volume whose name contains \DUroletitlereference{target\_cell\_} followed by at least one digit, followed by any other characters.
The following are all valid examples compatible with the default option: \DUroletitlereference{target\_cell\_5}, \DUroletitlereference{target\_cell\_2\_4}, \DUroletitlereference{target\_cell\_0\_0\_0}.
This can be altered by providing an optional second argument to the command invoking the macro.
Examples of alternate regular expressions to provide as the optional second argument include \DUroletitlereference{target\_cube}, \DUroletitlereference{scintillator\_volume{[}0-9{]}.?}, \DUroletitlereference{.*ej254.*}.
Please see {[}TRegexp{]}(\url{https://root.cern.ch/doc/master/classTRegexp.html}) for further details on how to construct a ROOT-compatible regular expression.

In this more general case, invoke the viewer as follows, where \DUroletitlereference{tcr} is the target-cell regex:

\texttt{`cpp
.L RATPACEventViewer2.cxx
RATPACEventViewer2(\textquotedbl{}some\_data\_file.root\textquotedbl{}, \textquotedbl{}tcr\textquotedbl{});
`}

\#\#\# Commands

Invoking the macro will automatically draw the detector geometry. The commands to then draw the particle tracks for a given event are as follows, where \DUroletitlereference{event} is the event number as defined in the original RAT-PAC TTree's 'mc.id':

\texttt{`cpp
drawTracks( <event> )
drawNextEvent()
drawPrevEvent()
`}

\#\#\# Output

In addition to drawing the particle tracks in the viewer, the macro will also print out a summary for each track, including particle name, number of points, and production and termination vertices.

\#\#\# Further Information

The comments in the macro file contain more-detailed information on usage. They are copied here below for convenience.


\section{Details%
  \label{details}%
}

\texttt{`cpp
// RATPACEventViewer2 -{}- for viewing RAT-PAC detector geometry and particle tracks in ROOT
//
// -{}- Usage: 'RATPACEventViewer2( <RAT-PAC rootfile>, {[}\textquotedbl{}target\_cell\_regex\textquotedbl{}{]} )' to draw detector;
// ~ ~ ~then 'drawTracks( <event number>, {[}IBD\_TF{]}, {[}highlight\_TF{]} )' to draw tracks for a given event
//
// -{}- Primarily written for IBD events, but should work just fine for anything
// ~ ~ ~with either one particle per top-level MC event (use IBD\_TF = kFALSE)
// ~ ~ ~or two (use IBD\_TF = kTRUE or simply omit IBD\_TF argument)
//
// -{}- \textquotedbl{}target\_cell\_regex\textquotedbl{} determines which detector volumes will be drawn;
// ~ ~ ~it defaults to \textquotedbl{}.*target\_cell\_{[}0-9{]}.*\textquotedbl{}, but any regex
// ~ ~ ~matching (exclusively) your desired volumes should work
// ~ ~ ~-{}- for more information on ROOT-compatible regexes, see:
// ~ ~ ~ ~ https://root.cern.ch/doc/master/classTRegexp.html
//
// -{}- \textquotedbl{}IBD\_TF\textquotedbl{} should be kTRUE for runs using the RAT-PAC IBD generator builtin
// ~ ~ ~and kFALSE otherwise
// ~ ~ ~-{}- note: in standard ROOT, kTRUE == true == 1 and kFALSE == false == 0
//
// -{}- \textquotedbl{}IBD\_TF\textquotedbl{} is kTRUE by default; this default can be changed in the function
// ~ ~ ~definition for \textquotedbl{}drawTracks()\textquotedbl{} if desired
//
// -{}- \textquotedbl{}highlight\_TF\textquotedbl{} simply switches cell highlighting on or off (default is on)
// ~ ~-{}- default for IBD events:
// ~ ~ ~ green = generation cell
// ~ ~ ~ red = positron-annihilation cell
// ~ ~ ~ blue = neutron-capture cell
//
// -{}- Example \textasciitilde{} Scan through some events:
// ~ ~ ~.L RATPACEventViewer2.cxx
// ~ ~ ~RATPACEventViewer2( \textquotedbl{}some\_data\_run.root\textquotedbl{} );
// ~ ~ ~drawTracks(0);
// ~ ~ ~drawNextEvent();
// ~ ~ ~drawNextEvent();
// ~ ~ ~drawNextEvent();
//
// -{}- As of this version, the following items are still hard-coded
// ~ ~and are tagged with //HC// in the code below:
// ~ ~- name of top \& world volumes
// ~ ~- shape of top \& world volumes
// ~ ~- dimensions of top \& world volumes
// ~ ~- shapes of target cells (currently \textquotedbl{}box\textquotedbl{})
// ~ ~- material of target cells (currently Eljen EJ-254 doped at 1.5\%wt Li-6)
//
// \textasciitilde{} Mark J. Duvall \textasciitilde{} mjduvall@hawaii.edu \textasciitilde{} Written 10/2019 \textasciitilde{} Updated 3/2020 \textasciitilde{} //
// \textasciitilde{} RATPACEventviewer v1.3.1 \textasciitilde{} //
`}

\end{document}
